\documentclass[11pt]{article}
\usepackage[a4paper, top=18mm, bottom=18mm, left=20mm, right=20mm]{geometry}
\usepackage[T2A]{fontenc}
\usepackage[utf8]{inputenc}
\usepackage[ukrainian]{babel}
\usepackage{graphicx}
\setlength{\parindent}{0pt}
\setlength{\parskip}{4pt}
\pagestyle{empty}
\begin{document}
%begin
{\large\textbf{closed-10}}

\textbf{Варіант \No\,1}\hfill \textit{ПІБ:}\,\underline{\hspace{60mm}}
\vspace{4mm}

%beginitem
1. Відмітити все, що є коректним оператором С++:
( 1 ) cout>>3;

( 2 ) int x+y=z;

( 3 ) cout<<3.4;

( 4 ) int f(int n);

%enditem
%beginitem
2. Відмітити все, що є коректним оператором С++:
( 1 ) x+3=5

( 2 ) n=2; n--

( 3 ) i,j=j,j

( 4 ) x=1//3

%enditem
%beginitem
3. Відмітити все, що є коректним оператором С++:
( 1 ) x+3=5

( 2 ) n=2; n--

( 3 ) i,j=j,j

( 4 ) x=1//3

%enditem
%beginitem
4. Відмітити все, що є коректним оператором С++:
( 1 ) x+3=5

( 2 ) n=2; n--

( 3 ) i,j=j,j

( 4 ) x=1//3

%enditem
%beginitem
5. Відмітити всі правильні варіанти:
( 1 ) 8,3

( 2 ) x*y

( 3 ) 0X2b

( 4 ) cpp

%enditem
%beginitem
6. Відмітити всі правильні варіанти:
( 1 ) 8,3

( 2 ) x*y

( 3 ) 0X2b

( 4 ) cpp

%enditem
%beginitem
7. Відмітити все, що є коректним оператором С++:
( 1 ) x+3=5

( 2 ) n=2; n--

( 3 ) i,j=j,j

( 4 ) x=1//3

%enditem
%beginitem
8. Відмітити все, що є коректним оператором С++:
( 1 ) cout>>3;

( 2 ) int x+y=z;

( 3 ) cout<<3.4;

( 4 ) int f(int n);

%enditem
%beginitem
9. Відмітити все, що є коректним оператором С++:
( 1 ) x+3=5

( 2 ) n=2; n--

( 3 ) i,j=j,j

( 4 ) x=1//3

%enditem
%beginitem
10. Відмітити все, що є коректним оператором С++:
( 1 ) x+3=5

( 2 ) n=2; n--

( 3 ) i,j=j,j

( 4 ) x=1//3

%enditem










%end
\newpage
\end{document}

